\pagenumbering{gobble}\section*{}
\vfill

{\bfseries Erläuterung zur Titelseite:}\\
Die Karte im Titelbild zeigt zwei Auschnitte aus der \textit{Uebersicht der im Königreich Hannover 1821–1844 von dem Herrn Hofrath Gauß gemessenen Dreieckssysteme erster Ordnung und derjenigen zweiter Ordnung die mit nahe gleicher Schärfe gemessen sind, zu Papen's topograph. Karte des Königreichs Hannover und Herzogthums Braunschweig, Originalmaßstab 1~:~1.000.000, Hannover 1845} als Teil des \textit{Topographischen Atlas des Königreichs Hannover und Herzogthums Braunschweig, Hannover 1832–1847.} Sie wird in der Kartenabteilung der Staatsbibliothek zu Berlin – Stiftung Preußischer Kulturbesitz aufbewahrt. (Kart.~N~25675)

Der bedeutende Mathematiker Carl Friedrich Gauß (1777–1855) war ab 1807 Professor an der Georg-August-Uni\-ver\-si\-tät und\linebreak Direktor der Sternwarte Göttingen. Seine trigonometrische Vermessung einiger neuer Gebiete des Königreichs Hannover bot die astronomisch-geodätische Grundlage für deren topographische Kartierung. Die gesamte Unternehmung ging als \glqq Gaußsche Landesaufnahme\grqq\  in die Geschichte ein.
\newpage
\pagenumbering{arabic}
%\setcounter{page}{1}
\section*{Inhalt}
\label{contents}
\newlength\contentspace
\setlength\contentspace{0.2em}

\vspace*{\contentspace}%
\noindent Fotos und Social Media \dotfill \pageref{Fotos und Social Media}

\vspace*{\contentspace}%
\noindent Exkursionen, Abendveranstaltung, Rahmenprogramm  \dotfill \pageref{exkursionen}

\vspace*{\contentspace}%
\noindent Platin- und Goldsponsoren \dotfill \pageref{platinsposoren}

\vspace*{\contentspace}%
\noindent Workshops am Mittwoch \dotfill \pageref{mittwoch-workshops}

\vspace*{\contentspace}%
\noindent Workshops am Donnerstag \dotfill \pageref{donnerstag-workshops}

\vspace*{\contentspace}%
\noindent Workshops am Freitag \dotfill \pageref{freitag-workshops}

\vspace*{\contentspace}%
\noindent Vorträge am Mittwoch \dotfill \pageref{mittwoch}

\vspace*{\contentspace}%
\noindent Vorträge am Donnerstag \dotfill \pageref{donnerstag}

\vspace*{\contentspace}%
\noindent Vorträge am Freitag \dotfill \pageref{freitag}

\vspace*{\contentspace}%
\noindent Posterausstellung \dotfill \pageref{donnerstag-poster}

\vspace*{\contentspace}%
\noindent OpenStreetMap-Samstag \dotfill \pageref{samstag}

\vspace*{\contentspace}%
\noindent Impressum, Medienpartner\dotfill \pageref{impressum}

%\vspace*{\contentspace}%
%\noindent Medienpartner \dotfill \pageref{medienpartner}
%
\vspace*{\contentspace}%
\noindent Campusplan und Raumpläne \dotfill \pageref{kartenseiten}

\justifying

\vspace{0.2cm}
\noindent
Die {\bfseries FOSSGIS-Konferenz 2026} wird vom gemeinnützigen FOSSGIS~e.V. und der OpenStreetMap-Community in Kooperation mit
dem Geographischen Institut der Georg-August-Universität Göttingen veranstaltet.  Die FOSSGIS-Konferenz ist eine Communityveranstaltung und wird vorwiegend ehrenamtlich organisiert. Ziel der Konferenz ist die Verbreitung von freier, quelloffener Software für Geoinformationssysteme. In den nächsten vier Tagen haben Sie
die Gelegenheit, sich mit Entwickler:innen und anderen Anwender:innen auszutauschen und neueste Informationen zu Anwendungen und Arbeitsmöglichkeiten zu erhalten.


\newpage
